% =============================================================================
%  Follow-up analyses carried out after the researcher review meeting.
%
%  Scope: only the additional work done in response to that feedback. The
%  original project, the dataset construction and the development of the
%  hidden Markov model are all assumed known and are not repeated here.
%
%  Figures are read from outputs/; regenerate with scripts/02 and scripts/08,
%  then compile twice with pdflatex.
% =============================================================================
\documentclass[11pt,a4paper]{article}

\usepackage[T1]{fontenc}
\usepackage[utf8]{inputenc}
\usepackage[margin=2.4cm]{geometry}
\usepackage{graphicx}
\usepackage{amsmath}
\usepackage{amssymb}
\usepackage{booktabs}
\usepackage{float}
\usepackage{caption}
\usepackage{enumitem}
\usepackage[hidelinks]{hyperref}

\graphicspath{
  {../outputs/patient_level/figures/}
  {../outputs/transitions/figures/}
  {../outputs/model/figures/}
  {figures/}
}

\captionsetup{font=small,labelfont=bf}
\setlist[itemize]{itemsep=2pt,topsep=4pt}
\setlist[enumerate]{itemsep=2pt,topsep=4pt}

\title{\textbf{Follow-up Analyses of Motor Subtype Progression}\\[0.6em]
       \large Additional work carried out after the research review meeting}
\author{Parv Patel \quad\textbullet\quad Mayank Tewatia\\[0.4em]
        \small Supervisor: Prof.\ K.~R.\ Sahasranand\\
        \small Department of Data Science, IIT Palakkad}
\date{2026}

\begin{document}
\maketitle

\begin{abstract}
\noindent
After reviewing our earlier work on TD/Indeterminate/PIGD motor subtype
progression in the PPMI cohort, several questions were raised about how we were
analysing and presenting the results. This report covers only the additional
analyses carried out in response, and assumes the earlier work is already
familiar.

Five things were asked for: a look at where inside the Indeterminate range those
observations actually fall; subtype progression shown for individual patients
rather than as a population summary; the conversion analysis repeated in the
direction of TD as well as PIGD; the raw measured tremor/PIGD ratio shown
directly rather than only through a fitted model; and the separate
population-level model plots combined into one.

The Indeterminate ratios were found to be concentrated near the two boundaries of
the Indeterminate range: about 78\% of the visits with a computable ratio occur
at just two values, one slightly above the PIGD threshold and one slightly below
the TD threshold. Correcting the patient-level progression analysis to use only
the 934 patients who have at least two labelled visits, 476 (51.0\%) show no
recorded subtype change, 210 (22.5\%) change once and 248 (26.6\%) change more
than once. Conversion is recorded in both directions, and requiring the new
subtype to persist to the next visit roughly halves the estimated rate in both
cases. The raw ratio shows a downward trend over follow-up, and also moves by
more than the entire width of the Indeterminate range between consecutive visits
of the same patient in about three quarters of cases.

The Indeterminate distribution is presented here as an observation, for
discussion, without further analysis of what follows from it. The remaining
results suggest that individual recorded subtype changes should be read
cautiously, and that a sustained change carries more weight than a single one.
\end{abstract}

% =============================================================================
\section{Purpose of the follow-up analysis}
% =============================================================================

The discussion of our earlier results raised a number of questions that we could
not answer from the analyses we had already run. They fell into five groups.

The first concerned the Indeterminate category. We had reported that it is
unstable, with only about a fifth of Indeterminate visits followed by another
Indeterminate visit. The question raised was where inside the Indeterminate range
those observations actually sit, and in particular whether they lie close to one
of the two subtype boundaries.

The second concerned how we were presenting progression. Our figures summarised
the whole population at each follow-up year, which mixes genuine change with the
fact that different patients are present at different times. What was wanted was
a view of what happens to individual patients.

The third was that our conversion analysis only looked in one direction. We had
estimated how long it takes to be recorded as PIGD, but not how often patients
move the other way.

The fourth concerned the ratio itself. The subtype label is produced by cutting a
continuous quantity into three categories, and our figure of ratio progression
showed a fitted model rather than the measurements. The request was to see the
raw values.

The fifth was a presentation issue: the model-based population results were
spread across three separate panels, which makes the three starting subtypes
difficult to compare.

We addressed each of these. None of the analyses below changes the dataset, the
subtype definition, the medication filter or any fitted model; they are new views
of data we already had, with one correction to a calculation that is described in
Section~\ref{sec:patient}.

For reference, the classification rule being used throughout is unchanged:
\[
  R = \frac{\text{mean of the 11 tremor items}}{\text{mean of the 5 gait and postural items}},
  \qquad
  \begin{cases}
    R \geq 1.15 & \text{TD} \\
    0.90 < R < 1.15 & \text{Indeterminate} \\
    R \leq 0.90 & \text{PIGD}
  \end{cases}
\]
In words: the ratio compares how severe a patient's tremor symptoms are against
how severe their gait and balance symptoms are, and two fixed cut-offs turn that
comparison into a label.

% =============================================================================
\section{Where inside the Indeterminate range do these observations sit?}
\label{sec:indeterminate}
% =============================================================================

The suggestion was to take the visits recorded as Indeterminate, particularly
those arriving from TD or from PIGD, and look at their actual ratio values rather
than only at how many there are. This section reports what the distribution looks
like. We have deliberately not taken the analysis further than that, so that the
observation can be discussed before anything is concluded from it.

There are 441 Indeterminate visits in the panel. Of these, 381 have a ratio that
can be computed; the remaining 60 have both subscores equal to zero, so the ratio
is undefined even though the classification rule still assigns them to
Indeterminate. Of the 381, 287 follow an earlier labelled visit and 94 are the
patient's first labelled visit, so they have no preceding subtype. Across the
whole panel, 332 transitions end in Indeterminate.

The observed ratios are concentrated near the two boundaries of the Indeterminate
range. Across all 381 visits only 11 distinct ratio values occur, and two of them
account for most of the observations:

\begin{itemize}
  \item \textbf{52.0\%} of these visits occur at $R = 0.9091$, which lies
        $0.009$ above the PIGD threshold of $0.90$;
  \item \textbf{26.2\%} occur at $R = 1.1364$, which lies $0.014$ below the TD
        threshold of $1.15$;
  \item together these two values account for \textbf{78.2\%} of the
        Indeterminate visits with a computable ratio.
\end{itemize}

The remaining nine values are spread thinly through the middle of the range, none
of them accounting for more than 11\% of observations.
Figure~\ref{fig:indeterminate-boundary} shows this distribution in its left
panel, with each bar split by the subtype recorded at the patient's previous
visit. Its right panel carries the accompanying next-visit breakdown from the
same analysis, which we are not discussing here.

\begin{figure}[H]
\centering
\includegraphics[width=\textwidth]{indeterminate_boundary.png}
\caption{Distribution of the tremor/PIGD ratio among visits recorded as
Indeterminate. \emph{Left:} every ratio value that occurs inside the range
$0.90 < R < 1.15$, across the 381 Indeterminate visits for which the ratio can be
computed. Bar length gives the number of visits and colour shows the subtype
recorded at the patient's previous visit, with visits that have no preceding
labelled visit shown separately in grey. The dashed lines mark the PIGD threshold
at $0.90$ and the TD threshold at $1.15$. \emph{Right:} for the 282 of those
visits that are followed by another labelled visit, the subtype recorded at that
next visit, grouped by where in the range the patient's ratio was sitting. The
figure was produced in response to the suggestion to examine where inside the
Indeterminate range these observations fall. The left panel shows a strong
concentration at the two values closest to the lower and upper boundaries; the
right panel is included as part of the analysis and is left for discussion.}
\label{fig:indeterminate-boundary}
\end{figure}

Broken down by where the patient came from, of the 287 Indeterminate visits that
follow an earlier labelled visit, 165 follow a visit recorded as TD, 75 follow a
visit recorded as PIGD and 47 follow another Indeterminate visit. The
concentration at the two boundary values is present in all three groups and in
the visits with no preceding label.

We are reporting this distribution as an observation and leaving its
interpretation open for discussion.

% =============================================================================
\section{Patient-level subtype progression}
\label{sec:patient}
% =============================================================================

Our earlier figures showed the share of the assessed population in each subtype
at each follow-up year. That summary has a known weakness: patients enter and
leave the study at different times, so a shift in the proportions mixes real
change with changes in who is being assessed. The request was to show individual
patients instead.

\subsection*{A correction to the population used}

Producing the patient-level version made us re-examine how we had been counting
subtype changes, and we found an error in the earlier figure.

There are 1{,}080 patients with at least one labelled visit, but 146 of them have
exactly one. A patient seen once cannot show a subtype change, because there is
no second observation to compare against. Our earlier calculation nevertheless
included them, and because their record collapses to a single subtype they were
counted among the patients who never changed. That inflated the apparent
stability of the cohort.

The correct population for any question about change is the \textbf{934 patients
with at least two labelled visits}. All patient-level results below use that
group, and the earlier figure has been regenerated.

\begin{table}[H]
\centering
\begin{tabular}{@{}lrr@{}}
\toprule
\textbf{Recorded subtype changes} & \textbf{Patients} & \textbf{Share} \\
\midrule
None                & 476 & \phantom{0}51.0\% \\
One                 & 210 & \phantom{0}22.5\% \\
Two or more         & 248 & \phantom{0}26.6\% \\
\midrule
Total               & 934 & 100.0\% \\
\bottomrule
\end{tabular}
\caption{Frequency of recorded subtype change among the 934 patients with at
least two labelled visits. The 146 patients with a single labelled visit are
excluded because they have no opportunity to show a change.}
\label{tab:change-frequency}
\end{table}

It is worth being explicit about what the first row means. ``No recorded
change'' means that the same subtype label was assigned at every labelled visit
the patient had. It does not mean the patient's condition was unchanged. A
patient can worsen substantially and still stay on the same side of the
threshold, and patients with short follow-up have fewer opportunities to be seen
changing. The figure is best read as a statement about the record rather than
about the disease.

\subsection*{Individual trajectories}

Figure~\ref{fig:patient-trajectories} shows every one of the 934 patients as a
single row. This is the direct answer to the request: nothing is averaged, and no
patient is left out.

\begin{figure}[H]
\centering
\includegraphics[width=0.78\textwidth]{patient_carpet.png}
\caption{Motor subtype trajectories for all 934 patients with at least two
labelled visits. Each horizontal row is one patient; colour shows the subtype
recorded at that patient's most recent visit, and blank space means follow-up had
ended. Rows are grouped by the subtype recorded at the first visit and then
ordered by how much of the follow-up period was recorded as PIGD. Reading down
each block shows how varied individual courses are, and the staircase shape on
the right is the loss of patients to follow-up, which is what made the earlier
population-level summary difficult to interpret.}
\label{fig:patient-trajectories}
\end{figure}

\subsection*{Recorded subtype courses}

A complementary way to summarise the same data is by the \emph{course} a patient
follows: the sequence of subtype labels across their visits, with repeats
collapsed, so that a patient recorded as TD, TD, PIGD, PIGD has the course
TD $\rightarrow$ PIGD. There are 86 distinct courses among the 934 patients. The
twelve commonest cover 775 patients (83.0\%); the remaining 159 patients are
spread across 74 rarer courses.

\begin{table}[H]
\centering
\begin{tabular}{@{}lrr@{}}
\toprule
\textbf{Recorded course} & \textbf{Patients} & \textbf{Share} \\
\midrule
TD                                                               & 399 & 42.7\% \\
TD $\rightarrow$ PIGD                                            & \phantom{0}83 & \phantom{0}8.9\% \\
PIGD                                                             & \phantom{0}71 & \phantom{0}7.6\% \\
TD $\rightarrow$ Indeterminate                                   & \phantom{0}39 & \phantom{0}4.2\% \\
TD $\rightarrow$ Indeterminate $\rightarrow$ TD                  & \phantom{0}35 & \phantom{0}3.7\% \\
Indeterminate $\rightarrow$ PIGD                                 & \phantom{0}27 & \phantom{0}2.9\% \\
Indeterminate $\rightarrow$ TD                                   & \phantom{0}26 & \phantom{0}2.8\% \\
PIGD $\rightarrow$ TD                                            & \phantom{0}24 & \phantom{0}2.6\% \\
TD $\rightarrow$ Indeterminate $\rightarrow$ PIGD                & \phantom{0}21 & \phantom{0}2.2\% \\
TD $\rightarrow$ PIGD $\rightarrow$ TD                           & \phantom{0}20 & \phantom{0}2.1\% \\
PIGD $\rightarrow$ TD $\rightarrow$ PIGD                         & \phantom{0}19 & \phantom{0}2.0\% \\
Indeterminate $\rightarrow$ TD $\rightarrow$ PIGD                & \phantom{0}11 & \phantom{0}1.2\% \\
\midrule
Twelve commonest courses                                         & 775 & 83.0\% \\
\bottomrule
\end{tabular}
\caption{The twelve commonest recorded subtype courses among the 934 patients
with at least two labelled visits. Courses ending where they started, such as
TD $\rightarrow$ Indeterminate $\rightarrow$ TD and PIGD $\rightarrow$ TD
$\rightarrow$ PIGD, together account for a noticeable share of patients.}
\label{tab:courses}
\end{table}

Figure~\ref{fig:subtype-courses} draws one representative patient for each of
these twelve courses, so that what a course looks like in time can be seen
directly. The representatives were picked as the patient with the typical number
of visits for that course rather than the longest followed, because choosing the
longest would make a decade of follow-up look normal when the median patient
contributes three visits over about two years.

\begin{figure}[H]
\centering
\includegraphics[width=0.82\textwidth]{patient_timelines.png}
\caption{One representative patient for each of the twelve commonest recorded
subtype courses. Dots mark actual visits and the coloured bar between two dots
shows the subtype recorded at the earlier of them. The number beside each row is
how many of the 934 patients follow that course, not how many patients are drawn;
one patient is drawn per course. Note the short follow-up in most rows, which is
typical of the cohort.}
\label{fig:subtype-courses}
\end{figure}

Two things stand out in Table~\ref{tab:courses}. Nearly half the patients have a
course consisting of a single subtype, which is consistent with
Table~\ref{tab:change-frequency}. And among those who do change, courses that
return to where they started are common: TD $\rightarrow$ Indeterminate
$\rightarrow$ TD, PIGD $\rightarrow$ TD $\rightarrow$ PIGD and TD $\rightarrow$
PIGD $\rightarrow$ TD together account for 74 patients. This is the same
instability we had already reported, seen from the level of the individual
patient.

% =============================================================================
\section{Conversion towards PIGD and towards TD}
% =============================================================================

Our earlier analysis estimated how long it takes for a patient to be recorded as
PIGD. The request was to run the same analysis in the opposite direction. We did
this using exactly the same method, so that the two are comparable: patients
count as being at risk if their first labelled visit is not the target subtype,
conversion is their first visit recorded as that subtype, patients who never
convert are censored at their last labelled visit, and a patient who moves back
and forth is counted only at their first qualifying entry.

Two definitions of conversion are reported in both directions. Under the
\emph{single-visit} definition, one visit recorded as the target subtype counts.
Under the \emph{sustained} definition, the target subtype must still be recorded
at the following visit. The gap between them measures how much of the apparent
conversion consists of one-off assessments that are not repeated.

\begin{table}[H]
\centering
\begin{tabular}{@{}llrrrr@{}}
\toprule
\textbf{Direction} & \textbf{Definition} & \textbf{At risk} &
\textbf{Conversions} & \textbf{5-year estimate} & \textbf{Median time} \\
\midrule
towards PIGD & single visit & 885 & 263 & 35.3\% & 3.00 years \\
towards PIGD & sustained    & 885 & 127 & 15.3\% & 3.00 years \\
\addlinespace
towards TD   & single visit & 304 & 125 & 61.2\% & 1.09 years \\
towards TD   & sustained    & 304 & \phantom{0}60 & 33.9\% & 1.33 years \\
\bottomrule
\end{tabular}
\caption{Recorded conversion in both directions, estimated by the same
survival method. The two at-risk groups are different populations and different
sizes, so the percentages are not a head-to-head comparison.}
\label{tab:conversion}
\end{table}

\begin{figure}[H]
\centering
\includegraphics[width=\textwidth]{conversion_both_directions.png}
\caption{Estimated cumulative probability of being recorded in the target subtype
over follow-up, for conversion towards PIGD (left) and towards TD (right), under
both definitions. Both panels use the same method, the same axes and the same
truncation point, so the shapes can be compared even though the two at-risk
groups differ. The dot marks the five-year estimate. Shaded bands are 95\%
intervals. The right-hand panel was added in response to the request to examine
conversion in the reverse direction.}
\label{fig:conversion-both}
\end{figure}

Three points about Table~\ref{tab:conversion} need stating carefully.

First, the two rows for each direction differ by a factor of roughly two.
Requiring a subtype to still be recorded at the next visit more than halves the
estimated five-year rate in both directions. Much of what looks like conversion
under the looser definition is a single assessment that is not repeated.

Second, the at-risk groups are not the same. The group at risk of being recorded
as PIGD is 885 patients; the group at risk of being recorded as TD is 304,
because most of the cohort starts as TD. The 61.2\% and 35.3\% figures are
therefore not directly comparable, and we would not want them presented side by
side as though they were.

Third, and most important, a recorded conversion towards TD should not be read as
clinical improvement. Under the single-visit definition, 61.2\% of the relevant
at-risk group were recorded as TD within five years, and half of those within
about one year. But the group at risk consists of patients whose first visit was
recorded as PIGD or Indeterminate, and the Indeterminate label in particular is
the least reliable of the three. A substantial part of this movement is likely to
be a first assessment that was not representative, corrected at the second visit,
rather than a patient whose symptoms genuinely shifted. The observed subtype can
also change again later, so a single TD observation is not a permanent state.

% =============================================================================
\section{The raw tremor/PIGD ratio over follow-up}
% =============================================================================

The subtype label is produced by cutting a continuous measurement into three
categories, so any statement about label stability is really a statement about
that measurement. Our earlier figure of ratio progression showed a fitted model,
and the request was to see the measurements themselves.
Figure~\ref{fig:raw-ratio-progression} does that: every observation is plotted,
each patient's visits are joined, and no model is involved.

\begin{figure}[H]
\centering
\includegraphics[width=0.90\textwidth]{raw_ratio_progression.png}
\caption{Raw observed tremor/PIGD ratio at every visit where it can be computed:
3{,}716 visits from 1{,}009 patients. Grey lines join visits belonging to the same
patient. The dashed lines are the TD threshold at $1.15$ and the PIGD threshold
at $0.90$, with the Indeterminate range shaded between them. The black line is
the median of the observations actually present in each year; it is a summary of
the raw data, not a fitted trend, and is affected by which patients are still
attending. The lower panel gives the number of visits contributing to each year.
No model output appears anywhere in this figure.}
\label{fig:raw-ratio-progression}
\end{figure}

Two things can be read from it. There is a downward trend: the median observed
ratio sits at about $1.82$ through the first three years and falls to about
$0.91$ by year 10. That is consistent with the direction we had estimated
previously, and it is visible in the measurements without any modelling. But the
later part of that line rests on very few patients --- by year 12 only 52 visits
remain --- and the median rises again there, which is a change in who is still
attending rather than a change in the patients. The count panel is part of the
reading.

The second and more important thing is how much scatter there is. The individual
patient lines cross the Indeterminate band repeatedly and in both directions. To
quantify this:

\begin{itemize}
  \item the standard deviation of the ratio \emph{within} a patient over time is
        about $1.22$, against an Indeterminate range only $0.25$ wide;
  \item about 54\% of the total variation in the ratio is between patients and
        46\% is within a patient over time, an intraclass correlation of
        $0.536$;
  \item the median change in a patient's own ratio from one visit to the next is
        $0.68$, and in 73\% of consecutive visit pairs the ratio moves by more
        than the entire width of the Indeterminate range.
\end{itemize}

Figure~\ref{fig:ratio-within-patient} shows the last of these directly.

\begin{figure}[H]
\centering
\includegraphics[width=0.72\textwidth]{ratio_within_patient.png}
\caption{Size of the change in a single patient's own tremor/PIGD ratio between
one visit and the next, across 2{,}707 consecutive visit pairs. The vertical line
marks the width of the whole Indeterminate range for comparison. Because this
compares the same patient measured twice, it isolates variation within a person
from variation between people.}
\label{fig:ratio-within-patient}
\end{figure}

The practical consequence is the one that matters for how the subtype labels
should be read. If the ratio routinely moves by more than the width of the band
between two visits of the same patient, then a patient whose ratio happens to sit
near a threshold will cross it from time to time for reasons unrelated to
progression. A single recorded change of subtype is therefore weak evidence on
its own, which is also what the gap between the two conversion definitions in
Table~\ref{tab:conversion} suggests.

% =============================================================================
\section{Combined population-level view of the model results}
% =============================================================================

The last request was a presentation one. The model-based population results had
been drawn as three separate panels, one for each starting subtype, which makes
it hard to compare them. Figure~\ref{fig:combined-hmm-population} puts all three
on one set of axes. This is the same model and the same numbers as before, drawn
differently; nothing about the model itself has changed.

\begin{figure}[H]
\centering
\includegraphics[width=\textwidth]{progression_combined.png}
\caption{Model-based probability of occupying each subtype over time, for
patients starting in each of the three subtypes, drawn on one set of axes.
Colour shows the subtype being predicted and line style shows the subtype the
patient starts in. Percentages are read at year 5. These are probabilities
produced by the fitted model for a patient's underlying subtype, not observed
proportions in the cohort. Replacing the earlier three-panel version allows the
three starting subtypes to be compared without moving between figures.}
\label{fig:combined-hmm-population}
\end{figure}

Read at five years, the model gives:

\begin{table}[H]
\centering
\begin{tabular}{@{}lrrr@{}}
\toprule
\textbf{Starting subtype} & \textbf{TD at 5 years} &
\textbf{Indeterminate at 5 years} & \textbf{PIGD at 5 years} \\
\midrule
TD            & 66\% & 15\% & 19\% \\
Indeterminate & 18\% & 19\% & 63\% \\
PIGD          & \phantom{0}4\% & 11\% & 86\% \\
\bottomrule
\end{tabular}
\caption{Model-based probabilities at a five-year horizon, by starting subtype.}
\label{tab:hmm-five-year}
\end{table}

Putting the three on one axis makes the asymmetry easy to see. A patient in TD
now has roughly a one-in-five chance of being in PIGD in five years, while a
patient in PIGD now has about a one-in-twenty-five chance of being in TD.
Patients currently in Indeterminate resemble PIGD far more than TD, with 63\% of
them in PIGD at five years. This is the same directional drift we see in the
observed transitions, expressed as a probability over a fixed horizon.

Two caveats attach to this figure. The probabilities refer to the model's
estimate of a patient's underlying subtype, not to the label that would be
recorded at a visit, and they are not observed proportions in the cohort. And the
projection assumes the fitted transition behaviour continues to apply over the
whole horizon; we have some evidence that rates are not constant over follow-up,
so the five-year values should be treated as indicative rather than exact.

% =============================================================================
\section{What these analyses taken together suggest}
% =============================================================================

The five pieces of work were requested separately, but they point in a
consistent direction.

About half of the patients who can be assessed show no recorded subtype change
at all during their follow-up, and the other half change at least once, often
more than once and often returning to where they started. Separately, the
Indeterminate observations were found to cluster at the two values closest to the
boundaries of the Indeterminate range. That distribution is recorded here as an
observation; what should be made of it is left for discussion.

The conversion analysis run in both directions shows that movement is recorded
towards TD as well as towards PIGD, and that requiring a change to persist to the
following visit roughly halves the estimated rate in both directions. The raw
ratio explains why: within a single patient the measurement moves by more than
the width of the Indeterminate range between most pairs of consecutive visits.

The practical conclusion is about interpretation rather than about the disease. A
single recorded change of subtype is not strong evidence that a patient's
condition has changed, and a change that persists across two visits should be
given considerably more weight than one that does not. Where a result depends on
counting subtype changes, it is worth checking how sensitive it is to that
distinction. This applies to our own earlier results as much as to anything else,
and Table~\ref{tab:conversion} is the clearest illustration of the size of the
effect.

Separately, correcting the population used for the patient-level analysis matters
in its own right. Restricting to patients who can actually demonstrate a change
lowered the proportion recorded as never changing, and we have used the corrected
figure of 51.0\% throughout.

% =============================================================================
\section{Current status and the next validation step}
% =============================================================================

All five of the requested analyses are complete. The distribution of ratios
within the Indeterminate range was examined and is reported in
Section~\ref{sec:indeterminate} for discussion. The patient-level progression
analysis was
corrected and regenerated on the appropriate population of 934 patients.
Conversion towards TD was added using the same method as the existing PIGD
analysis. The raw ratio was examined directly, without a model. The
population-level model results were combined into a single figure.

A stricter patient-level leave-one-out validation of the existing hidden Markov
model is currently being performed as a separate analysis. In that analysis each
patient is held out completely, the model is fitted using the remaining patients,
and the held-out patient's next recorded subtype is predicted one visit at a time
using only the history available before each prediction. Those results will be
reported separately.

\vspace{1em}
\noindent
\small
Data used in the preparation of this report were obtained from the Parkinson's
Progression Markers Initiative (PPMI) database.

\end{document}
